\section{模型假设}

\begin{enumerate}[label=\textbf{A\arabic*.}]
  \item 将附件中00:10对应数据视为区间$[0{:}00,0{:}10)$的平均功率，第144个值对应$[23{:}50,24{:}00)$；若官方澄清为区间起点，所有序列整体一致修正。
  \item 充电效率与放电效率均取0.9，充放电量均在交流母线侧计量；总往返效率解释为0.9的口径留作敏感性比较。
  \item 不允许向外网售电；光伏或已购电配额过剩时允许形成无收益的富余电量。
  \item 问题一日初、日末储能量均为6000 kWh。问题二至四跨日连续传递，不进行每日重置。
  \item 问题二在每日0:00只能使用目标日以前的数据。日前购电和名义电池轨迹固定，日内偏差仅由紧急购电和富余处置补救。
  \item 问题三主结算按最终执行购电量相对0:00原计划一次结算；下调电量按0.5倍电价退回。下调另加罚的解释作为备选口径。
  \item 题面没有给出并网购电上限、售电价、固定调整费和电池退化费，主模型不人为增加这些参数。
\end{enumerate}

\section{符号说明}

一天划分为$T=144$个时段，$\Delta t=1/6$ h。主要符号见表\ref{tab:symbols}。

\begin{table}[H]
\centering
\caption{主要符号及单位}
\label{tab:symbols}
\begin{tabularx}{\textwidth}{>{\centering\arraybackslash}p{2.4cm}X>{\centering\arraybackslash}p{2.6cm}}
\toprule
符号 & 含义 & 单位 \\
\midrule
$L_t,G_t$ & 第$t$时段负荷电量、光伏电量 & kWh \\
$p_t$ & 第$t$时段外购电价 & 元/kWh \\
$q_t,r_t,e_t$ & 日前计划、调整后正常购电、紧急购电量 & kWh \\
$c_t,d_t$ & 交流母线侧充电量、放电量 & kWh \\
$s_t$ & 第$t$时段开始时的电池内部储能量 & kWh \\
$u_t$ & 未被负荷或电池利用的富余电量 & kWh \\
$\eta_c,\eta_d$ & 充、放电效率，主模型均取0.9 & 无量纲 \\
$\Omega,\pi_\omega$ & 场景集合、场景概率 & -- \\
\bottomrule
\end{tabularx}
\end{table}

